Transformation optics employs coordinate transformations to design materials with
various interesting electromagnetic properties. Perhaps the most famous
application is invisibility cloaking. In 2006, Pendry et al. described the
parameters necessary to construct an invisibility cloak ADD REFERENCE. They
expressed this description through the matrix:
\begin{equation*}
\boldsymbol{\varepsilon}
=\boldsymbol{\mu}=
\begin{bmatrix}
\frac{r-a}{r} & 0 & 0\\
0 & \frac{r}{r-a} & 0\\
0 & 0 & \left(\frac{b}{b-a}\right)^2 \frac{r-a}{r}
\end{bmatrix}.
\end{equation*}
What does this matrix represent? Why does it produce invisibility? What
assumptions underlie this result? Under what circumstances does it work?

This document aims to formulate these questions more rigorously, answer
them, and provide intuition for understanding the results.
